\section{The Method of Undetermined Coefficients I}

\noindent
In this section we consider the constant coefficient equation
\begin{equation} \label{eq:5.4.1}
ay''+by'+cy=e^{\alpha x}G(x),
\end{equation}
where $\alpha$ is a constant and $G$ is a polynomial.

From Theorem~\ref{thmtype:5.3.2}, the general solution of (\ref{eq:5.4.1})
is
$y=y_p+c_1y_1+c_2y_2$, where $y_p$ is a particular solution of
(\ref{eq:5.4.1}) and $\{y_1,y_2\}$ is a fundamental set of
solutions of the complementary equation
$$
ay''+by'+cy=0.
$$
In Section~5.2 we showed how to find $\{y_1,y_2\}$. In this
section we'll show how to find $y_p$. The procedure that we'll use
is called {\color{blue}\it the method of undetermined coefficients\/}.

Our first example is similar to
Exercises~\ref{exer:5.3.16}--\ref{exer:5.3.21}.

\begin{example}\label{example:5.4.1} \rm
Find a particular solution of
\begin{equation} \label{eq:5.4.2}
y''-7y'+12y=4e^{2x}.
\end{equation}
Then find the general solution.
\end{example}

\solution
Substituting $y_p=Ae^{2x}$ for $y$ in (\ref{eq:5.4.2}) will produce a
constant multiple of $Ae^{2x}$ on the left side of (\ref{eq:5.4.2}), so it
may be possible to choose $A$ so that $y_p$ is a solution of
(\ref{eq:5.4.2}). Let's try it;   if $y_p=Ae^{2x}$ then
$$
y_p''-7y_p'+12y_p=4Ae^{2x}-14Ae^{2x}+12Ae^{2x}=2Ae^{2x}=4e^{2x}
$$
if $A=2$. Therefore $y_p=2e^{2x}$ is a particular solution of
(\ref{eq:5.4.2}). To find the general solution, we note that the
characteristic polynomial of the complementary equation
\begin{equation} \label{eq:5.4.3}
y''-7y'+12y=0
\end{equation}
is $p(r)=r^2-7r+12=(r-3)(r-4)$, so $\{e^{3x},e^{4x}\}$ is a
fundamental set of solutions of (\ref{eq:5.4.3}). Therefore the general
solution of (\ref{eq:5.4.2}) is
$$
 y=2e^{2x}+c_1e^{3x}+c_2e^{4x}.
$$

\begin{example}\label{example:5.4.2} \rm
Find a particular solution of
\begin{equation} \label{eq:5.4.4}
y''-7y'+12y=5e^{4x}.
\end{equation}
Then find the general solution.
\end{example}

\solution
Fresh from our success in finding a particular solution of
(\ref{eq:5.4.2}) --- where we chose $y_p=Ae^{2x}$ because the right side
of
(\ref{eq:5.4.2}) is a constant multiple of $e^{2x}$ --- it may seem
reasonable to try $y_p=Ae^{4x}$ as a particular solution of
(\ref{eq:5.4.4}). However, this won't work, since we saw in
Example~\ref{example:5.4.1} that $e^{4x}$ is a solution of the
complementary equation (\ref{eq:5.4.3}), so substituting $y_p=Ae^{4x}$
into the left side of (\ref{eq:5.4.4}) produces zero on the left, no
matter how we choose$A$. To discover a suitable form for $y_p$,
we use the same approach that we used in Section~5.2 to find a
second solution of
$$
ay''+by'+cy=0
$$
in the case where the characteristic equation has a repeated real
root: we look for solutions of (\ref{eq:5.4.4}) in the form $y=ue^{4x}$,
where $u$ is a function to be determined. Substituting
\begin{equation} \label{eq:5.4.5}
y=ue^{4x},\quad
y'=u'e^{4x}+4ue^{4x},\mbox{\quad and \quad}
y''=u''e^{4x}+8u'e^{4x}+16ue^{4x}
\end{equation}
into (\ref{eq:5.4.4}) and canceling the common factor $e^{4x}$ yields
$$
(u''+8u'+16u)-7(u'+4u)+12u=5,
$$
or
$$
u''+u'=5.
$$
By inspection we  see that $u_p=5x$ is a particular solution of
this equation, so $y_p=5xe^{4x}$ is a particular solution of
(\ref{eq:5.4.4}). Therefore
$$
y=5xe^{4x}+c_1e^{3x}+c_2e^{4x}
$$
is the general solution.

\begin{example}\label{example:5.4.3} \rm
Find a particular solution of
\begin{equation} \label{eq:5.4.6}
y''-8y'+16y=2e^{4x}.
\end{equation}
\end{example}

\solution
Since the characteristic polynomial of the complementary equation
\begin{equation} \label{eq:5.4.7}
y''-8y'+16y=0
\end{equation}
is $p(r)=r^2-8r+16=(r-4)^2$, both $y_1=e^{4x}$ and $y_2=xe^{4x}$ are
solutions of (\ref{eq:5.4.7}). Therefore (\ref{eq:5.4.6}) does not have a
solution of the form $y_p=Ae^{4x}$ or $y_p=Axe^{4x}$. As in
Example~\ref{example:5.4.2}, we look for solutions of (\ref{eq:5.4.6}) in the
form $y=ue^{4x}$, where $u$ is a function to be determined.
Substituting from (\ref{eq:5.4.5}) into (\ref{eq:5.4.6}) and canceling the
common factor $e^{4x}$ yields
$$
(u''+8u'+16u)-8(u'+4u)+16u=2,
$$
or
$$
u''=2.
$$
Integrating twice and taking the constants of integration to be zero
shows that $u_p=x^2$ is a particular solution of this equation, so
$y_p=x^2e^{4x}$ is a particular solution of (\ref{eq:5.4.4}). Therefore
$$
y=e^{4x}(x^2+c_1+c_2x)
$$
is the general solution. \bbox

The preceding examples illustrate the following facts concerning
the form of a particular solution $y_p$
of a constant coefficent equation
$$
ay''+by'+cy=ke^{\alpha x},
$$
where $k$ is a nonzero constant:

\medskip

\begin{alist}
\item % (a)
If $e^{\alpha x}$ isn't  a solution of the complementary
equation
\begin{equation} \label{eq:5.4.8}
ay''+by'+cy=0,
\end{equation}
then  $y_p=Ae^{\alpha x}$, where $A$ is a constant. (See
Example~\ref{example:5.4.1}).
\item % (b)
If $e^{\alpha x}$ is a solution of (\ref{eq:5.4.8}) but $xe^{\alpha x}$
is not, then  $y_p=Axe^{\alpha x}$,  where $A$ is a constant.
(See Example~\ref{example:5.4.2}.)
\item % (c)
If both $e^{\alpha x}$ and $xe^{\alpha x}$ are solutions of (\ref{eq:5.4.8}),
then  $y_p=Ax^2e^{\alpha x}$,  where $A$ is a constant.
(See Example~\ref{example:5.4.3}.)
\end{alist}
See Exercise~\ref{exer:5.4.30}  for the proofs of these facts.

In all three cases you can just substitute the appropriate form for
$y_p$ and its derivatives directly into
$$
ay_p''+by_p'+cy_p=ke^{\alpha x},
$$
and solve for the constant $A$, as we did in
Example~\ref{example:5.4.1}.
 (See Exercises~\ref{exer:5.4.31}--\ref{exer:5.4.33}.)
However, if the equation is
$$
ay''+by'+cy=k e^{\alpha x}G(x),
$$
where $G$ is a polynomial of degree greater than zero, we recommend
that you use the substitution $y=ue^{\alpha x}$ as we did in
Examples~\ref{example:5.4.2} and \ref{example:5.4.3}. The equation for $u$
will turn out to be
\begin{equation} \label{eq:5.4.9}
au''+p'(\alpha)u'+p(\alpha)u=G(x),
\end{equation}
where $p(r)=ar^2+br+c$ is the characteristic polynomial of the
complementary equation and $p'(r)=2ar+b$ (Exercise~\ref{exer:5.4.30});
however, you shouldn't memorize this since it's easy to derive the
equation for $u$ in any particular case. Note, however, that if
$e^{\alpha x}$ is a solution of the complementary equation then
$p(\alpha)=0$, so (\ref{eq:5.4.9}) reduces to
$$
au''+p'(\alpha)u'=G(x),
$$
while if both $e^{\alpha x}$ and $xe^{\alpha x}$ are solutions of the
complementary equation then $p(r)=a(r-\alpha)^2$ and
$p'(r)=2a(r-\alpha)$, so $p(\alpha)=p'(\alpha)=0$ and (\ref{eq:5.4.9})
reduces to
$$
au''=G(x).
$$

\begin{example}\label{example:5.4.4} \rm
Find a particular solution of
\begin{equation} \label{eq:5.4.10}
y''-3y'+2y=e^{3x}(-1+2x+x^2).
\end{equation}
\end{example}

\solution
Substituting
$$
y=ue^{3x},\quad y'=u'e^{3x}+3ue^{3x},\mbox{\quad and }
y''=u''e^{3x}+6u'e^{3x}+9ue^{3x}
$$
into (\ref{eq:5.4.10}) and  canceling $e^{3x}$  yields
$$
(u''+6u'+9u)-3(u'+3u)+2u=-1+2x+x^2,
$$
or
\begin{equation} \label{eq:5.4.11}
u''+3u'+2u=-1+2x+x^2.
\end{equation}
As in Example~\ref{exer:5.3.2}, in order to guess a form for a particular
solution of
(\ref{eq:5.4.11}), we note that substituting a second degree polynomial
$u_p=A+Bx+Cx^2$ for $u$ in the left side of (\ref{eq:5.4.11}) produces
another second degree polynomial with coefficients that depend upon
$A$, $B$, and $C$;   thus,
$$
\mbox{ if \quad}u_p=A+Bx+Cx^2\mbox{\quad then \quad}
u_p'=B+2Cx\mbox{\quad and \quad} u_p''=2C.
$$
If $u_p$ is to satisfy (\ref{eq:5.4.11}),  we must have
\begin{eqnarray*}
u_p''+3u_p'+2u_p&=&2C+3(B+2Cx)+2(A+Bx+Cx^2)\\
&=&(2C+3B+2A)+(6C+2B)x+2Cx^2=-1+2x+x^2.
\end{eqnarray*}
Equating  coefficients of like powers of $x$ on the two sides of the
last equality yields
$$
\begin{array}{rcr}
2C&=&1\phantom{.}\\
2B+6C&=&2\phantom{.}\\
2A+3B+2C&=& -1.
\end{array}
$$
Solving these equations for $C$, $B$, and $A$ (in that order) yields
 $C=1/2,B=-1/2,A=-1/4$.  Therefore
$$
u_p=-{1\over4}(1+2x-2x^2)
$$
is a particular solution of  (\ref{eq:5.4.11}), and
$$
y_p=u_pe^{3x}=-{e^{3x}\over4}(1+2x-2x^2)
$$
is a particular solution of  (\ref{eq:5.4.10}).

\begin{example}\label{example:5.4.5} \rm
 Find a particular solution of
\begin{equation} \label{eq:5.4.12}
y''-4y'+3y=e^{3x}(6+8x+12x^2).
\end{equation}
\end{example}

\solution
Substituting
$$
y=ue^{3x},\quad y'=u'e^{3x}+3ue^{3x},\mbox{\quad and }
y''=u''e^{3x}+6u'e^{3x}+9ue^{3x}
$$
into (\ref{eq:5.4.12}) and  canceling $e^{3x}$  yields
$$
(u''+6u'+9u)-4(u'+3u)+3u=6+8x+12x^2,
$$
or
\begin{equation} \label{eq:5.4.13}
u''+2u'=6+8x+12x^2.
\end{equation}
There's no $u$ term in this equation, since $e^{3x}$ is a solution of
the complementary equation for (\ref{eq:5.4.12}). (See
Exercise~\ref{exer:5.4.30}.) Therefore (\ref{eq:5.4.13}) does not have a
particular solution of the form $u_p=A+Bx+Cx^2$ that we used
successfully in Example~\ref{example:5.4.4}, since with this choice of
$u_p$,
$$
u_p''+2u_p'=2C+(B+2Cx)
$$
can't contain the last term ($12x^2$) on the right side of
(\ref{eq:5.4.13}). Instead, let's try $u_p=Ax+Bx^2+Cx^3$ on the grounds
that
$$
u_p'=A+2Bx+3Cx^2\mbox{\quad and \quad}
u_p''=2B+6Cx
$$
together contain all the powers of $x$ that appear on the right side
of (\ref{eq:5.4.13}).

Substituting these expressions in place of $u'$ and $u''$  in
(\ref{eq:5.4.13}) yields
$$
(2B+6Cx)+2(A+2Bx+3Cx^2)=(2B+2A)+(6C+4B)x+6Cx^2=6+8x+12x^2.
$$
Comparing coefficients of like powers of $x$ on the two sides of the
last equality shows that $u_p$ satisfies (\ref{eq:5.4.13}) if
$$
\begin{array}{rcr}
6C&=&12\phantom{.}\\
4B+6C&=&8\phantom{.}\\
2A+2B\phantom{+6u_2}&=&6.
\end{array}
$$
Solving these equations successively yields $C=2$, $B=-1$, and $A=4$.
Therefore
$$
u_p=x(4-x+2x^2)
$$
is a particular solution of  (\ref{eq:5.4.13}),
and
$$
y_p=u_pe^{3x}=xe^{3x}(4-x+2x^2)
$$
is a particular solution of  (\ref{eq:5.4.12}).

\begin{example}\label{example:5.4.6} \rm
Find a particular solution of
\begin{equation} \label{eq:5.4.14}
4y''+4y'+y=e^{-x/2}(-8+48x+144x^2).
\end{equation}
\end{example}

\solution  Substituting
$$
y=ue^{-x/2},\quad  y'=u'e^{-x/2}-{1\over2}ue^{-x/2},\mbox{\quad and
\quad} y''=u''e^{-x/2}-u'e^{-x/2}+{1\over4}ue^{-x/2}
$$
into (\ref{eq:5.4.14}) and  canceling $e^{-x/2}$  yields
$$
4\left(u''-u'+{u\over4}\right)+4\left(u'-{u\over2}\right)+u=4u''=-8+48x+144x^2,
$$
or
\begin{equation} \label{eq:5.4.15}
u''=-2+12x+36x^2,
\end{equation}
which does not contain $u$ or $u'$ because $e^{-x/2}$ and $xe^{-x/2}$
are both solutions of the complementary equation. (See
Exercise~\ref{exer:5.4.30}.) To obtain a particular solution of
(\ref{eq:5.4.15}) we integrate twice, taking the constants of integration
to be zero;   thus,
$$
u_p'=-2x+6x^2+12x^3\mbox{\quad and \quad}
u_p=-x^2+2x^3+3x^4=x^2(-1+2x+3x^2).
$$
Therefore
$$
y_p=u_pe^{-x/2}=x^2e^{-x/2}(-1+2x+3x^2)
$$
is a particular solution of  (\ref{eq:5.4.14}).

\boxit{Summary}

\noindent
The preceding examples illustrate the following facts concerning
 particular solutions of a constant coefficent equation of the form
$$
ay''+by'+cy=e^{\alpha x}G(x),
$$
where $G$ is a polynomial (see Exercise~\ref{exer:5.4.30}):
\begin{alist}
\item % (a)
If $e^{\alpha x}$ isn't  a solution of the complementary
equation
\begin{equation} \label{eq:5.4.16}
ay''+by'+cy=0,
\end{equation}
then  $y_p=e^{\alpha x}Q(x)$,
 where $Q$ is a polynomial of the same degree as $G$.
(See Example~\ref{example:5.4.4}).
\item % (b)
If $e^{\alpha x}$ is a solution of (\ref{eq:5.4.16}) but $xe^{\alpha x}$
is not, then  $y_p=xe^{\alpha x}Q(x)$,
 where $Q$ is a polynomial of the same degree as $G$.
(See Example~\ref{example:5.4.5}.)
\item % (c)
If both $e^{\alpha x}$ and $xe^{\alpha x}$ are solutions of (\ref{eq:5.4.16}),
then  $y_p=x^2e^{\alpha x}Q(x)$,
 where $Q$ is a polynomial of the same degree as $G$.
(See Example~\ref{example:5.4.6}.)
\end{alist}

In all three cases, you can just substitute the appropriate form for
$y_p$ and its derivatives directly into
$$
ay_p''+by_p'+cy_p=e^{\alpha x}G(x),
$$
and solve for the coefficients of the polynomial $Q$. However, if you
try this you will see that the computations are more tedious than
those that you encounter by making the substitution $y=ue^{\alpha x}$
and finding a particular solution of the resulting equation for $u$.
(See Exercises~\ref{exer:5.4.34}-\ref{exer:5.4.36}.) In Case~\part{a}
the equation for $u$ will be of the form
$$
au''+p'(\alpha)u'+p(\alpha)u=G(x),
$$
with a particular solution of the form $u_p=Q(x)$, a polynomial of the
same degree as $G$, whose coefficients can be found by the method used
in Example~\ref{example:5.4.4}. In Case~\part{b}  the equation for
$u$ will be of the form
$$
au''+p'(\alpha)u'=G(x)
$$
(no $u$ term on the left), with a particular solution of the form
$u_p=xQ(x)$, where $Q$ is a polynomial of the same degree as $G$ whose
coefficents can be found by the method used in
Example~\ref{example:5.4.5}. In Case~\part{c}  the equation for $u$
will be of the form
$$
au''=G(x)
$$
with a particular solution of the form $u_p=x^2Q(x)$ that can be
obtained by integrating $G(x)/a$ twice and taking the constants of
integration to be zero, as in Example~\ref{example:5.4.6}.

\boxit{Using the Principle of Superposition}

\noindent
The next example shows how to combine the method of undetermined
coefficients and Theorem~\ref{thmtype:5.3.3}, the principle of
superposition.

\begin{example}\label{example:5.4.7} \rm
Find a particular solution of
\begin{equation} \label{eq:5.4.17}
y''-7y'+12y=4e^{2x}+5e^{4x}.
\end{equation}
\end{example}

\solution
In Example~\ref{example:5.4.1} we found that $y_{p_1}=2e^{2x}$
is a particular solution of
$$
y''-7y'+12y=4e^{2x},
$$
and in Example~\ref{example:5.4.2} we found that $y_{p_2}=5xe^{4x}$
is a particular solution of
$$
y''-7y'+12y=5e^{4x}.
$$
Therefore the principle of superposition implies that
$y_p=2e^{2x}+5xe^{4x}$ is a particular solution of (\ref{eq:5.4.17}).

\exercises
In Exercises \ref{exer:5.4.1}--\ref{exer:5.4.14} find a particular
solution.

\begin{exerciselist}

\begin{tabular}[t]{@{}p{168pt}@{}p{168pt}}
\item\label{exer:5.4.1} $y''-3y'+2y=e^{3x}(1+x)$
&\vspace*{10pt}  \item\label{exer:5.4.2} $y''-6y'+5y=e^{-3x}(35-8x)$
\end{tabular}

\begin{tabular}[t]{@{}p{168pt}@{}p{168pt}}
\item\label{exer:5.4.3} $y''-2y'-3y=e^x(-8+3x)$
&  \item\label{exer:5.4.4} $y''+2y'+y=e^{2x}(-7-15x+9x^2)$
\end{tabular}

\begin{tabular}[t]{@{}p{168pt}@{}p{168pt}}
\item\label{exer:5.4.5} $y''+4y=e^{-x}(7-4x+5x^2)$
&  \item\label{exer:5.4.6} $y''-y'-2y=e^x(9+2x-4x^2)$
\end{tabular}

\begin{tabular}[t]{@{}p{168pt}@{}p{168pt}}
\item\label{exer:5.4.7} $y''-4y'-5y=-6xe^{-x}$
&  \item\label{exer:5.4.8} $y''-3y'+2y=e^x(3-4x)$
\end{tabular}

\begin{tabular}[t]{@{}p{168pt}@{}p{168pt}}
\item\label{exer:5.4.9} $y''+y'-12y=e^{3x}(-6+7x)$
&  \item\label{exer:5.4.10} $2y''-3y'-2y=e^{2x}(-6+10x)$
\end{tabular}

\begin{tabular}[t]{@{}p{168pt}@{}p{168pt}}
\item\label{exer:5.4.11} $y''+2y'+y=e^{-x}(2+3x)$ &
\item\label{exer:5.4.12} $y''-2y'+y=e^x(1-6x)$
\end{tabular}

\item\label{exer:5.4.13} $y''-4y'+4y=e^{2x}(1-3x+6x^2)$

\item\label{exer:5.4.14} $9y''+6y'+y=e^{-x/3}(2-4x+4x^2)$

\exercisetext{In Exercises \ref{exer:5.4.15}--\ref{exer:5.4.19} find the general
solution.}

\begin{tabular}[t]{@{}p{168pt}@{}p{168pt}}
\item\label{exer:5.4.15} $y''-3y'+2y=e^{3x}(1+x)$ &
\vspace*{10pt}\item\label{exer:5.4.16} $y''-6y'+8y=e^x(11-6x)$
\end{tabular}

\begin{tabular}[t]{@{}p{168pt}@{}p{168pt}}
\item\label{exer:5.4.17} $ y''+6y'+9y=e^{2x}(3-5x)$ &
\item\label{exer:5.4.18} $y''+2y'-3y=-16xe^x$
\end{tabular}

\item\label{exer:5.4.19} $y''-2y'+y=e^x(2-12x)$

\newpage
\exercisetext{In Exercises \ref{exer:5.4.20}--\ref{exer:5.4.23} solve the
initial value problem and plot the solution.}

\item\label{exer:5.4.20} \CGex $y''-4y'-5y=9e^{2x}(1+x), \quad  y(0)=0,\quad y'(0)=-10$
\item\label{exer:5.4.21} \CGex $y''+3y'-4y=e^{2x}(7+6x), \quad  y(0)=2,\quad y'(0)=8$
\item\label{exer:5.4.22} \CGex $y''+4y'+3y=-e^{-x}(2+8x), \quad  y(0)=1,\quad y'(0)=2$
\item\label{exer:5.4.23} \CGex $y''-3y'-10y=7e^{-2x}, \quad  y(0)=1,\quad y'(0)=-17$

\exercisetext{In Exercises~\ref{exer:5.4.24}--\ref{exer:5.4.29} use
the principle of superposition to find a particular solution.}

\item\label{exer:5.4.24} $y''+y'+y=xe^x+e^{-x}(1+2x)$
\item\label{exer:5.4.25} $y''-7y'+12y=-e^x(17-42x)-e^{3x}$
\item\label{exer:5.4.26} $y''-8y'+16y=6xe^{4x}+2+16x+16x^2$
\item\label{exer:5.4.27} $y''-3y'+2y=-e^{2x}(3+4x)-e^x$
\item\label{exer:5.4.28} $y''-2y'+2y=e^x(1+x)+e^{-x}(2-8x+5x^2)$
\item\label{exer:5.4.29}
$y''+y=e^{-x}(2-4x+2x^2)+e^{3x}(8-12x-10x^2)$

\item\label{exer:5.4.30}
\begin{alist}
 \item % (a)
Prove that  $y$ is a solution of the constant coefficient equation
$$
ay''+by'+cy=e^{\alpha x}G(x)
\eqno{\rm (A)}
$$
if and only if $y=ue^{\alpha x}$, where $u$ satisfies
$$
au''+p'(\alpha)u'+p(\alpha)u=G(x)
\eqno{\rm (B)}
$$
and   $p(r)=ar^2+br+c$ is the characteristic polynomial of
the complementary equation
$$
ay''+by'+cy=0.
$$

For the rest of this exercise, let $G$ be
 a polynomial. Give the requested proofs for the case where
$$
G(x)=g_0+g_1x+g_2x^2+g_3x^3.
$$

\item % (b)
Prove that if $e^{\alpha x}$ isn't  a solution of the complementary
equation then (B) has a particular solution of the
form $u_p=A(x)$, where $A$ is a polynomial of the same degree as $G$,
as in Example~\ref{example:5.4.4}. Conclude that (A) has
a particular solution of the form $y_p=e^{\alpha x}A(x)$.

\item % (c)
Show that if $e^{\alpha x}$ is a solution of the complementary
equation and $xe^{\alpha x}$ isn't , then (B) has a
particular solution of the form $u_p=xA(x)$, where $A$ is a polynomial
of the same degree as $G$, as in Example~\ref{example:5.4.5}. Conclude
that (A) has a particular solution of the form
$y_p=xe^{\alpha x}A(x)$.

\item % (d)
Show that if $e^{\alpha x}$ and $xe^{\alpha x}$ are both solutions of
the complementary equation then (B) has a
particular solution of the form $u_p=x^2A(x)$, where $A$ is a
polynomial of the same degree as $G$, and $x^2A(x)$ can be obtained by
integrating $G/a$ twice, taking the constants of integration to be
zero, as in Example~\ref{example:5.4.6}. Conclude that
(A) has a particular solution of the form
$y_p=x^2e^{\alpha x}A(x)$.
\end{alist}

\exercisetext{Exercises~\ref{exer:5.4.31}--\ref{exer:5.4.36}
treat the equations considered in
Examples~\ref{example:5.4.1}--\ref{example:5.4.6}. Substitute the suggested
form of $y_p$ into the equation and equate the resulting coefficients
of like functions on the two sides of the resulting equation to derive
a set of simultaneous equations for the coefficients in $y_p$. Then
solve for the coefficients to obtain $y_p$. Compare the work you've
done with the work required to obtain the same results in
Examples~\ref{example:5.4.1}--\ref{example:5.4.6}.}

\item\label{exer:5.4.31}
Compare with Example~\ref{example:5.4.1}:
$$
y''-7y'+12y=4e^{2x};\quad y_p=Ae^{2x}
$$

\item\label{exer:5.4.32}
Compare with Example~\ref{example:5.4.2}:
$$
y''-7y'+12y=5e^{4x};\quad y_p=Axe^{4x}
$$

\item\label{exer:5.4.33}
Compare with Example~\ref{example:5.4.3}.
$$
y''-8y'+16y=2e^{4x};\quad y_p=Ax^2e^{4x}
$$

\item\label{exer:5.4.34}
Compare with Example~\ref{example:5.4.4}:
$$
y''-3y'+2y=e^{3x}(-1+2x+x^2),\quad y_p=e^{3x}(A+Bx+Cx^2)
$$

\item\label{exer:5.4.35}
Compare with Example~\ref{example:5.4.5}:
$$
y''-4y'+3y=e^{3x}(6+8x+12x^2),\quad y_p=e^{3x}(Ax+Bx^2+Cx^3)
$$

\item\label{exer:5.4.36}
Compare with Example~\ref{example:5.4.6}:
$$
4y''+4y'+y=e^{-x/2}(-8+48x+144x^2),\quad y_p=e^{-x/2}(Ax^2+Bx^3+Cx^4)
$$

\item\label{exer:5.4.37}
Write $y=ue^{\alpha x}$ to find the general
solution.

\begin{tabular}[t]{@{}p{168pt}@{}p{168pt}}
\part{a} $y''+2y'+y=\dst{e^{-x}\over\sqrt x}$ &
\part{b} $y''+6y'+9y=e^{-3x}\ln x$\\
\part{c} $y''-4y'+4y=\dst{e^{2x}\over1+x}$ &
\part{d} $4y''+4y'+y=\dst{4e^{-x/2}\left({1\over x}+x\right)}$
\end{tabular}

\item\label{exer:5.4.38}
Suppose $\alpha\ne0$ and $k$ is a positive integer. In most
calculus books integrals like $\int x^k e^{\alpha x}\,dx$ are
evaluated by integrating by parts $k$ times. This exercise presents
another method. Let
$$
y=\int e^{\alpha x}P(x)\,dx
$$
with
$$
P(x)=p_0+p_1x+\cdots+p_kx^k, \mbox{\quad (where $p_k\ne0$)}.
$$
\begin{alist}
\item %  (a)
Show that $y=e^{\alpha x}u$, where
$$
u'+\alpha u=P(x).
\eqno{\rm (A)}
$$
\item % (b)
Show that  (A) has a particular solution of the form
$$
u_p=A_0+A_1x+\cdots+A_kx^k,
$$
where $A_k$, $A_{k-1}$, \dots, $A_0$ can be computed successively
by equating coefficients of $x^k,x^{k-1}, \dots,1$ on both sides of the
equation
$$
u_p'+\alpha u_p=P(x).
$$
\item % (c)
Conclude that
$$
\int e^{\alpha x}P(x)\,dx=\left(A_0+A_1x+\cdots+A_kx^k\right)e^{\alpha x}
+c,
$$
where $c$ is a constant of integration.
\end{alist}

\item\label{exer:5.4.39}
Use the method of Exercise~\ref{exer:5.4.38} to evaluate the integral.

\begin{tabular}[t]{@{}p{168pt}@{}p{168pt}}
 \part{a} $  \int e^x(4+x)\,dx$
&  \part{b} $\int e^{-x}(-1+x^2)\,dx$\\[10pt]
 \part{c} $\int x^3e^{-2x}\,dx$
&  \part{d} $\int e^x(1+x)^2\,dx$\\[10pt]
 \part{e} $\int e^{3x}(-14+30x+27x^2)\,dx$
&  \part{f} $\int e^{-x}(1+6x^2-14x^3+3x^4)\,dx$
\end{tabular}

\item\label{exer:5.4.40}
Use the method suggested in Exercise~\ref{exer:5.4.38} to
evaluate $\int x^ke^{\alpha x}\,dx$, where $k$ is an arbitrary
positive integer and $\alpha\ne0$.
\end{exerciselist}
